\documentclass[11pt,a4paper]{article}

\usepackage[utf8]{inputenc}
\usepackage[T1]{fontenc}
\usepackage{lmodern}
\usepackage[slovak]{babel}
\usepackage[a4paper,margin=2.2cm]{geometry}
\usepackage[hidelinks]{hyperref}
\usepackage[protrusion=true,expansion=false]{microtype}
\usepackage{parskip}

% Spoločné údaje projektu — pri zmene upravte iba tento súbor.
\newcommand{\projectTitleSK}{Detekcia falošných online recenzií na základe jazykových a behaviorálnych príznakov}
\newcommand{\projectTitleEN}{Detection of Fake Online Reviews Based on Linguistic and Behavioral Features}
\newcommand{\projectAcronym}{FAKERADAR}
\newcommand{\projectId}{Zatiaľ nepridelené (študentský návrh)}
\newcommand{\projectDate}{01.11.2026}
\newcommand{\projectStart}{01/2027}
\newcommand{\projectEnd}{12/2029}

\newcommand{\applicant}{Slovenská technická univerzita v Bratislave, Fakulta informatiky a informačných technológií}
\newcommand{\applicantShort}{FIIT STU}
\newcommand{\applicantAddress}{Ilkovičova 6276/2, 842 16 Bratislava 4}
\newcommand{\applicantICO}{00397687}
\newcommand{\applicantDIC}{2020845255}
\newcommand{\applicantICDPH}{SK2020845255}
\newcommand{\applicantEmail}{info@fiit.stuba.sk}
\newcommand{\applicantPhone}{+421 2 210 44 111}
\newcommand{\statutory}{prof. Ing. Ivan Kotuliak, PhD.}
\newcommand{\statutoryRole}{dekan}
\newcommand{\piName}{Ing. Ivan Kapustík}
\newcommand{\piEmail}{ivan.kapustik@stuba.sk} % e-mail potvrdený v oficiálnych materiáloch FIIT

\newcommand{\budgetAPVV}{268\,000}
\newcommand{\budgetCoFin}{32\,000}
\newcommand{\budgetTotal}{300\,000}

% Anotácia = presné znenie schváleného zamerania (odovzdané 04.10.2026, hodnotené cvičiacim 5/5 b);
% zdroj: old_zameranie_clanok.tex (abstrakt).
\newcommand{\projectAnnotation}{%
Online recenzie výrazne ovplyvňujú rozhodovanie spotrebiteľov, no ich dôveryhodnosť narúšajú falošné recenzie generované firmami, platenými recenzentmi či automatizovanými nástrojmi. Tento článok prináša prehľad súčasných metód detekcie falošných recenzií so zameraním na porovnanie jazykových (textových) a behaviorálnych (na správaní recenzenta založených) príznakov.}


\setlength{\parindent}{0pt}
\setlength{\emergencystretch}{3em}

\begin{document}

\begin{center}
{\Large\bfseries Zameranie projektu}\\[0.4em]
{\large Názov projektu, akronym a anotácia}\\[0.8em]
{\small Schválené zameranie — odovzdané 04.10.2026, hodnotené cvičiacim na maximum (5/5 b).}\\
{\small Pokrýva položky VV-A1-03 (názov), VV-A1-04 (akronym) a VV-A1-09 (anotácia, max. 2\,000 znakov) formulára „Vzor žiadosti pre aplikovaný výskum“ (APVV VV 2024, Príloha 5).}
\end{center}

\vspace{0.8em}

\section*{VV-A1-03 — Názov projektu}
\textbf{Názov v slovenskom jazyku:} \projectTitleSK

\textbf{Názov v anglickom jazyku:} \projectTitleEN\\
{\small\emph{Anglický názov nebol v schválenom zameraní uvedený; pre dvojjazyčný formulár používame tento preklad.}}

\section*{VV-A1-04 — Akronym projektu}
\projectAcronym\\
{\small\emph{Akronym nebol v schválenom zameraní uvedený; pre potreby formulára VV-A1-04 navrhujeme tento akronym.}}

\section*{VV-A1-09 — Anotácia projektu (max. 2\,000 znakov)}
\projectAnnotation

\vspace{0.6em}
Anotácia má 369 znakov vrátane medzier (limit 2\,000 znakov).

\vspace{1.5em}
{\footnotesize Poznámka: slovenský názov a anotácia zodpovedajú presnému zneniu schváleného zamerania (súbor \nolinkurl{old_zameranie_clanok.tex}, odovzdaný 04.10.2026). Anglický názov a akronym sú odvodené pre potreby formulárov; ich finálne znenie potvrdí vyučujúci.}

\end{document}
